\section{Structured Expressivity Certificates}
\label{sec:structured-expressivity}

The full-SPD channel can realize every \(R\in\SPD^m\).  Implementable
optimizers often restrict \(P\) to a coordinate, block, Kronecker, or other
structured family.  For a family \(\Fcal\subset\SPD^d\) equipped with a
within-family distance \(d_\Fcal\), define the visible restricted complexity
\begin{equation}
  D_{\Fcal}^{\mathrm{vis}}(A,R;P_0)
  =
  \inf_{\substack{P\in\Fcal\\A^\top PA=R}}d_\Fcal(P_0,P),
  \label{eq:visible-restricted-complexity}
\end{equation}
with value \(+\infty\) when no structured completion exists.  This section
characterizes finiteness for diagonal and block families.

\subsection{Diagonal geometry}

Let \(a_j^\top\) be the \(j\)-th row of \(A\in\R^{d\times m}\) and set
\(G_j=a_ja_j^\top\in\PSD^m\).  A positive diagonal cometric
\(D(p)=\diag(p_1,\ldots,p_d)\) has visible image
\begin{equation}
  A^\top D(p)A=\sum_{j=1}^d p_jG_j.
  \label{eq:diagonal-image}
\end{equation}
Let \(\Kcal_{\mathrm{diag}}=\cone(G_1,\ldots,G_d)\).

\begin{theorem}[Diagonal expressivity and facial alternative]
\label{thm:diagonal-expressivity}
\mbox{}\par\noindent
Let \(D_0=\diag(p_1^0,\ldots,p_d^0)\) with \(p_j^0>0\), and define
\begin{equation}
  D_{\mathrm{diag}}^{\mathrm{vis}}
  =
  \inf_{\substack{p_j>0\\\sum_jp_jG_j=R}}
  \left(\sum_{j=1}^d\log^2\frac{p_j}{p_j^0}\right)^{1/2}.
  \label{eq:diagonal-complexity}
\end{equation}
Then \(D_{\mathrm{diag}}^{\mathrm{vis}}<\infty\) if and only if
\begin{equation}
  R\in\ri\Kcal_{\mathrm{diag}},
  \label{eq:diagonal-ri}
\end{equation}
equivalently, \(R=\sum_jp_jG_j\) for some \(p_j>0\).  When finite, the
infimum is attained.  If the \(G_j\) are linearly independent, the feasible
diagonal cometric is unique.

There are two distinct failure certificates.
\begin{enumerate}[label=(\roman*),leftmargin=*]
  \item If \(R\notin\Kcal_{\mathrm{diag}}\), there exists symmetric \(Y\)
  such that
  \[
    \tr(YR)<0,
    \qquad
    \tr(YG_j)\ge0\quad\forall j.
  \]
  \item If
  \(R\in\Kcal_{\mathrm{diag}}\setminus\ri\Kcal_{\mathrm{diag}}\), there
  exists symmetric \(Y\) such that
  \[
    \tr(YR)=0,
    \qquad
    \tr(YG_j)\ge0\quad\forall j,
    \qquad
    \tr(YG_{j_\star})>0\quad\text{for some }j_\star.
  \]
\end{enumerate}
The first strictly separates an infeasible target; the second exposes the
proper face containing a boundary target that requires some diagonal weights
to vanish.
\end{theorem}

\begin{example}[An end-to-end visible training channel]
\label{ex:worked-training-channel}
Let the rows of \(A\in\R^{3\times2}\) be
\[
  a_1^\top=(1,0),
  \qquad
  a_2^\top=(0,1),
  \qquad
  a_3^\top=(1,1).
\]
Take \(P_0=I_3\) and the visible target \(R=I_2\).  Since
\[
  R_0=A^\top A=\begin{pmatrix}2&1\\1&2\end{pmatrix},
\]
the full-SPD canonical completion is
\[
  P_\star=\Psi_{I_3,A}(I_2)
  =\frac19
  \begin{pmatrix}
    8&-1&-2\\
    -1&8&-2\\
    -2&-2&5
  \end{pmatrix},
  \qquad A^\top P_\star A=I_2.
\]
Thus the visible request has a unique nearest full-SPD realization.  Within
the diagonal family, however, it requires the semidefinite coefficient vector
\(p=(1,1,0)\); no strictly positive diagonal geometry realizes it.  The
matrix
\[
  Y=\begin{pmatrix}0&1\\1&0\end{pmatrix}
\]
is a facial certificate: \(\tr(YR)=0\),
\(\tr(YG_1)=\tr(YG_2)=0\), and \(\tr(YG_3)=2\), so it is a nontrivial
functional on \(\operatorname{span}\Kcal_{\mathrm{diag}}\).

For the visible covector \(\alpha=Ae_1=(1,0,1)^\top\), the canonical update
is
\[
  u_\star=-P_\star\alpha
  =\left(-\frac23,\frac13,-\frac13\right)^\top,
  \qquad A^\top u_\star=-e_1.
\]
Every feasible full geometry has this same visible update, while its invisible
component may differ by any vector in
\(\ker A^\top=\operatorname{span}\{(1,1,-1)^\top\}\).  The example
therefore displays, in one calculation, exact visible reduction, canonical
full completion, invisible drift, and a structured-family boundary
certificate.
\end{example}

\subsection{Block geometry}

Partition the rows of \(A\) into blocks \(A_b\).  A block cometric has the
form \(P=\blockdiag(P_1,\ldots,P_B)\), \(P_b\succ0\), and visible image
\begin{equation}
  A^\top PA=\sum_bA_b^\top P_bA_b.
  \label{eq:block-image}
\end{equation}
Define
\[
  \Kcal_{\mathrm{block}}
  =\left\{\sum_bA_b^\top P_bA_b:P_b\succeq0\right\}.
\]

\begin{theorem}[Block expressivity]
\label{thm:block-expressivity}
\mbox{}\par\noindent
For a block reference \(P_0=\blockdiag(P_1^0,\ldots,P_B^0)\), the quantity
\begin{equation}
  D_{\mathrm{block}}^{\mathrm{vis}}
  =
  \inf_{\substack{P_b\succ0\\\sum_bA_b^\top P_bA_b=R}}
  \left(\sum_b\dai(P_b^0,P_b)^2\right)^{1/2}
  \label{eq:block-complexity}
\end{equation}
is finite if and only if
\(R\in\ri\Kcal_{\mathrm{block}}\), equivalently if positive definite
blocks realize \(R\).  When finite, the infimum is attained.

The cone \(\Kcal_{\mathrm{block}}\) is closed.  If
\(R\notin\Kcal_{\mathrm{block}}\), there is symmetric \(Y\)
such that
\[
  \tr(YR)<0,
  \qquad
  A_bYA_b^\top\succeq0\quad\forall b.
\]
If
\(R\in\Kcal_{\mathrm{block}}\setminus
\ri\Kcal_{\mathrm{block}}\), there is a facial certificate with
\[
  \tr(YR)=0,
  \qquad
  A_bYA_b^\top\succeq0\quad\forall b,
  \qquad
  A_{b_\star}YA_{b_\star}^\top\ne0\quad\text{for some }b_\star.
\]
\end{theorem}

The proofs in \cref{app:expressivity-proofs} use the identity that a linear
map sends the relative interior of a finite-dimensional convex cone onto the
relative interior of its image.  The distinction between strict separation
and facial exposure is essential: targets on a proper face are realizable by
semidefinite structured geometries while remaining unreachable by strictly
positive optimizer blocks.  In both facial alternatives, nontriviality is
measured on the linear span of the image cone; ambient nonzero \(Y\) alone is
insufficient because it may annihilate the entire cone.

The certificates have direct finite-dimensional feasibility forms.  For the
diagonal boundary case, scaling permits the normalization
\(\sum_j\tr(YG_j)=1\), leaving linear equalities and inequalities in \(Y\);
strict infeasibility can be normalized by \(\tr(YR)\le-1\).  For block
geometry, the corresponding constraints
\(A_bYA_b^\top\succeq0\) together with
\(\sum_b\tr(A_bYA_b^\top)=1\) form an SDP feasibility problem.  These
formulations make the certificates computable in principle; algorithmic
complexity is outside the present characterization.  In finite precision,
the normalization equality, cone residual, and strict separation margin
should be reported together: a small residual with a vanishing margin signals
proximity to a proper face, whereas the exact alternatives in the theorems
concern exact feasibility.

Kronecker, low-rank, and tensor families require separate image geometry and
are not covered by the diagonal/block theorems.
